\chapter{Topology in Lean}\label{ch:topology_in_lean}

% Chapter orientation paragraph
In this chapter, we examine how topology is formalized in Lean and
Mathlib, building on the type class machinery introduced in the
previous chapters.
We begin by discussing how real numbers are represented in Mathlib
and why many operations on them are noncomputable
(Section~\ref{sec:reals_and_topspaces}).
We then introduce Mathlib's
\lstinline[language=lean]|TopologicalSpace| type class and the
definition of continuity in terms of open sets.
In Section~\ref{sec:filters}, we discuss
filters---Mathlib's central tool for reasoning uniformly about
limits, convergence, and continuity---and show how they connect
global and local notions of continuity.
Finally, in Section~\ref{sec:connectivity}, we recall the
definitions of connectedness and path-connectedness that will be
needed in the next chapter.

\section{Real Numbers and Topological
  Spaces}\label{sec:reals_and_topspaces}

% REVISED: P42 — removed computable reals red herring entirely.
In the previous chapters, we built rational numbers from integers
using quotient types.
In Mathlib, real numbers are constructed by a similar process: they
are defined as the \textbf{Cauchy completion} of $\mathbb{Q}$,
that is, equivalence classes of Cauchy sequences of rational
numbers~\cite{mathlibdoc}.
Unlike our construction of \lstinline[language=lean]|myRat|, where
classical axioms were used only for convenience (to avoid writing a
decision procedure for inequality), the standard
\lstinline[language=lean]|Real| type relies on classical axioms more
fundamentally.
Many operations on real numbers---such as the sine function, the
exponential, or division---are defined using the axiom of choice and
do not correspond to executable algorithms.
Consequently, they are marked as
\lstinline[language=lean]|noncomputable|:
\begin{lstlisting}[language=lean]
noncomputable def realSin (x : ℝ) : ℝ := Real.sin x
\end{lstlisting}
This marker indicates that Lean's kernel cannot evaluate the
function by computation, but we can still state and prove theorems
about it.

Real numbers in Mathlib are registered as instances of numerous
algebraic and topological type classes, including: linear ordered
field, normed space, metric space, and topological
space~(\cite{mathlib2020}).
Since $\mathbb{R}$ is a metric space, notions like convergence can
be defined using the familiar $\varepsilon$-$\delta$ formulation:
\begin{lstlisting}[language=lean]
def ConvergesTo (s : ℕ → ℝ) (a : ℝ) :=
  ∀ ε > 0, ∃ N, ∀ n ≥ N, |s n - a| < ε
\end{lstlisting}
where \lstinline[language=lean]|s| is a sequence and
\lstinline[language=lean]|a| is the limit point.
More generally, these concepts are formalized using topology.

\subsection{The \lstinline[language=lean]|TopologicalSpace| type
  class}

In standard textbooks such as Munkres~\cite{munkres}, a
topological space is defined as a set $X$ equipped with a
collection $\mathcal{T}$ of subsets of $X$, called \textbf{open
  sets}, satisfying the following axioms:
\begin{itemize}
  \item The whole space $X$ and the empty set $\varnothing$ are
        open.
  \item The intersection of any two open sets is open.
  \item The union of any collection of open sets is open.
\end{itemize}

% REVISED: P43 — fixed TopologicalSpace definition with proper types
This is reflected in Mathlib's
\lstinline[language=lean]|TopologicalSpace| type class:
\begin{lstlisting}[language=lean]
class TopologicalSpace (α : Type*) where
  IsOpen : Set α → Prop
  isOpen_univ : IsOpen Set.univ
  isOpen_inter : ∀ (s t : Set α), IsOpen s → IsOpen t → IsOpen (s ∩ t)
  isOpen_sUnion : ∀ (S : Set (Set α)), (∀ t ∈ S, IsOpen t) → IsOpen (sUnion S)
\end{lstlisting}
Recall that Lean treats sets as predicates
(\lstinline[language=lean]|Set α := α → Prop|), where set
membership is function application.
Here, \lstinline[language=lean]|IsOpen| is a predicate on sets
indicating whether they are open.
The three fields correspond to the standard topological axioms:
\begin{itemize}
  \item The whole space is open
        (\lstinline[language=lean]|isOpen_univ|).
  \item The intersection of two open sets is open
        (\lstinline[language=lean]|isOpen_inter|).
  \item The union of any collection of open sets is open
        (\lstinline[language=lean]|isOpen_sUnion|).
\end{itemize}
\begin{note}
  The empty set is open, as it is the union of the empty collection
  of open sets:
  $\varnothing = \bigcup \varnothing$.
\end{note}

Using this type class, global continuity can be defined in terms of
open sets, and indeed this is how it is defined in Mathlib:
\begin{lstlisting}[language=lean]
structure Continuous (f : X → Y) : Prop where
  isOpen_preimage : ∀ s, IsOpen s → IsOpen (f ⁻¹' s)
\end{lstlisting}
A function $f : X \to Y$ between topological spaces is continuous
if the preimage of every open set in $Y$ is open in $X$.
This is equivalent to the familiar $\varepsilon$-$\delta$ definition
when $X$ and $Y$ are metric spaces.

However, when proving results about limits, convergence, and
continuity, Mathlib predominantly uses an alternative framework
based on \textbf{filters}.
This approach may seem more abstract initially, but it provides a
more general and powerful framework that works uniformly across all
topological spaces and, by extension, metric spaces.

% Section transition
We now introduce filters and explain why Mathlib adopts them as the
primary language for reasoning about convergence.

\section{Filters}\label{sec:filters}

% REVISED: P44 — fixed all typos throughout this section
The concept of a limit is quite broad: there are many types of
limits to consider.
For instance, the limit of a sequence or a function at a point,
limits at infinity (from above or below), one-sided limits (from the
left or right), and so on.
Defining each of these separately would require a huge amount of
work in Mathlib, with significant duplication of theorems and
proofs.
Moreover, many fundamental theorems (like the characterization of
continuity via limits) would need to be reproved for each type of
limit.
Fortunately, Bourbaki solved this problem by introducing the notion
of \textbf{filters}, which unify all concepts of limits,
convergence, neighborhoods, and terms like ``eventually'' or
``frequently'' into a single framework.

Intuitively, a filter represents a notion of ``sufficiently large''
subsets.
More formally, a filter $F$ on a type $X$ is a collection of
subsets of $X$ satisfying three axioms:
\begin{lstlisting}[language=lean]
structure Filter (α : Type*) where
  sets : Set (Set α)
  univ_sets : Set.univ ∈ sets
  sets_of_superset {x y : Set α} : x ∈ sets → x ⊆ y → y ∈ sets
  inter_sets {x y : Set α} : x ∈ sets → y ∈ sets → x ∩ y ∈ sets
\end{lstlisting}
The field \lstinline[language=lean]|sets| suggests thinking of a
filter as a collection of sets.
The three axioms correspond to:
\begin{itemize}
  \item $X \in F$ (the whole space is in the filter).
  \item If $U \in F$ and $U \subseteq V$, then $V \in F$
        (supersets of ``large'' sets are ``large'').
  \item If $U, V \in F$, then $U \cap V \in F$
        (finite intersections of ``large'' sets are ``large'').
\end{itemize}
Note that if a filter $F$ contains the empty set, then it contains
all subsets of $X$ (since every set is a superset of
$\varnothing$).
Such a filter is called the \textbf{bottom filter} and is
degenerate; non-trivial filters never contain the empty set.

\subsection{Key filter concepts}

Filters are quite categorical objects, and we will use them
intuitively rather than developing the full theory.
In particular, we will use the following concepts:

\begin{itemize}
  \item \textbf{Neighborhood filter}
        \lstinline[language=lean]|𝓝 x|:
        In a topological space, this filter contains all
        neighborhoods of the point $x$.
        A set is in \lstinline[language=lean]|𝓝 x| if it contains
        an open set containing $x$.
        This captures the idea of ``near $x$.''

  \item \textbf{At-top filter}
        \lstinline[language=lean]|atTop : Filter ℕ|:
        Contains sets that include all sufficiently large natural
        numbers.
        Formally, $U \in$
        \lstinline[language=lean]|atTop| if and only if there
        exists $N$ such that
        $\{n \mid n \geq N\} \subseteq U$.
        This captures the idea of ``$n \to \infty$.''

  \item \lstinline[language=lean]|∀ᶠ x in f, p x|
        (\lstinline[language=lean]|f.Eventually p|):
        ``Eventually in filter $f$, property $p$ holds.''
        This means there exists some set $U \in f$ such that $p$
        holds for all $x \in U$.

  \item \lstinline[language=lean]|∃ᶠ x in f, p x|
        (\lstinline[language=lean]|f.Frequently p|):
        ``Frequently in filter $f$, property $p$ holds.''
        This means for every set $U \in f$, there exists some
        $x \in U$ where $p$ holds.
        This captures the idea that $p$ holds ``infinitely often''
        or ``arbitrarily close.''

  \item \lstinline[language=lean]|Tendsto f l₁ l₂|:
        ``Function $f$ tends from filter $l_1$ to filter $l_2$.''
        This is used for convergence.
\end{itemize}

% ADDED: P45 — the actual definition of Tendsto, which was missing
\subsection{The \lstinline[language=lean]|Tendsto| predicate}

The predicate \lstinline[language=lean]|Tendsto| is the key
definition that makes filters useful for convergence.
Its definition is remarkably concise:
\begin{lstlisting}[language=lean]
def Tendsto (f : α → β) (l₁ : Filter α) (l₂ : Filter β) : Prop :=
  l₁.map f ≤ l₂
\end{lstlisting}
where \lstinline[language=lean]|Filter.map f l₁| is the
\textbf{pushforward} of the filter $l_1$ along $f$: a set $U$
belongs to \lstinline[language=lean]|l₁.map f| if and only if
$f^{-1}(U) \in l_1$.
The ordering $\leq$ on filters is reverse inclusion of their set of
members: $F \leq G$ means that every set in $G$ is also in $F$.
Unwinding the definitions,
\lstinline[language=lean]|Tendsto f l₁ l₂| is equivalent to:

$$
  \forall\, U \in l_2,\quad f^{-1}(U) \in l_1.
$$

In other words, the preimage of every ``large'' set in $l_2$ is
``large'' in $l_1$.

\begin{example}
  We can express convergence of a sequence $s_n$ to a limit $a$
  using filters:
  \lstinline[language=lean]|Tendsto s atTop (𝓝 a)| means
  $s_n \to a$ as $n \to \infty$.

  Unwinding the definition: for every neighborhood $U$ of $a$
  (i.e., every $U \in \mathcal{N}(a)$), the preimage
  $s^{-1}(U) = \{n \in \mathbb{N} \mid s_n \in U\}$ is in
  \lstinline[language=lean]|atTop|, meaning it contains all
  sufficiently large $n$.
  This recovers the familiar definition: for every open set $U$
  containing $a$, there exists $N$ such that $s_n \in U$ for all
  $n \geq N$.
\end{example}

\subsection{From global to local continuity}

The definition of \lstinline[language=lean]|Tendsto| also provides
a natural bridge between global continuity (defined in terms of
open sets) and local continuity (defined at individual points).
As seen before, the structure
\lstinline[language=lean]|Continuous| defines global continuity
in terms of preimages of open sets.
Mathlib also provides alternative definitions for local
continuity:
\lstinline[language=lean]|ContinuousAt| defines continuity at a
single point,
\lstinline[language=lean]|ContinuousWithinAt| defines continuity
within a set at a point, and
\lstinline[language=lean]|ContinuousOn| defines continuity on an
entire set.
All these local characterizations are defined in terms of filters.

The connection between global and local continuity is established
by the following fundamental theorem:
\begin{lstlisting}[language=lean]
theorem Continuous.tendsto (hf : Continuous f) (x) :
    Tendsto f (𝓝 x) (𝓝 (f x)) :=
  ((nhds_basis_opens x).tendsto_iff <| nhds_basis_opens <| f x).2
    fun t ⟨hxt, ht⟩ => ⟨f ⁻¹' t, ⟨hxt, ht.preimage hf⟩, Subset.rfl⟩
\end{lstlisting}
The statement
\lstinline[language=lean]|Tendsto f (𝓝 x) (𝓝 (f x))| is a
reformulation of continuity at $x$ in terms of neighborhood
filters: for every neighborhood $U$ of $f(x)$, the preimage
$f^{-1}(U)$ is a neighborhood of $x$.

The proof uses \textbf{filter bases}.
Similar to how a topological space can be defined via a basis of
open sets, a filter can be defined via a basis of sets.
A basis $B$ for a filter $F$ is a nonempty collection of sets
that is closed under finite intersections (up to refinement):
for any two sets $U, V \in B$, there exists a set $W \in B$
such that $W \subseteq U \cap V$.
\begin{lstlisting}[language=lean]
structure FilterBasis (α : Type*) where
  sets : Set (Set α)
  nonempty : sets.Nonempty
  inter_sets {x y : Set α} : x ∈ sets → y ∈ sets → ∃ z ∈ sets, z ⊆ x ∩ y
\end{lstlisting}
Given a basis, we can generate a filter and prove properties
about it by working only with basis elements.

\newpage
To better understand the proof, we can expand it using tactic
mode:
\begin{lstlisting}[language=lean]
theorem Continuous'.tendsto (hf : Continuous f) (x) :
    Tendsto f (𝓝 x) (𝓝 (f x)) := by
  rw [(nhds_basis_opens x).tendsto_iff (nhds_basis_opens (f x))]
  intro t ⟨hft_in, ht_open⟩
  use f ⁻¹' t
  constructor
  · exact ⟨hft_in, ht_open.preimage hf⟩
  · exact Subset.rfl
\end{lstlisting}
The lemma \lstinline[language=lean]|nhds_basis_opens x| states
that the neighborhood filter \lstinline[language=lean]|𝓝 x| has
a basis consisting of all open sets containing $x$.
The rewrite
\lstinline[language=lean]|rw [(nhds_basis_opens x).tendsto_iff (nhds_basis_opens (f x))]|
reduces the goal to: for every open neighborhood $t$ of $f(x)$,
there exists an open neighborhood $s$ of $x$ with
$f(s) \subseteq t$.
We introduce an arbitrary open neighborhood $t$ of $f(x)$ using
\lstinline[language=lean]|intro t ⟨hft_in, ht_open⟩|.
The natural witness is the preimage
\lstinline[language=lean]|f ⁻¹' t|.
The first subgoal---that $x \in f^{-1}(t)$ and $f^{-1}(t)$ is
open---follows from \lstinline[language=lean]|hft_in| and from
\lstinline[language=lean]|ht_open.preimage hf|, which uses the
global continuity of $f$ to conclude that the preimage of an
open set is open.
The second subgoal---that
$\forall\, x \in f^{-1}(t),\, f(x) \in t$---holds by definition
of preimage.

Here is the full example:
[\href{https://live.lean-lang.org/\#codez=JYWwDg9gTgLgBAWQIYwBYBtgCMBQEwCmAdnACr4ToQDmAnnAGLDowFRwDKB8A8lACZscOQQDM4kAM4B9KASTo4ALgC8nbgDoAkhGBwAFAAZlcQLiEAShxEIRAMYRwAVxhIs6AnDFxJwIgWm2jlAAbh6qcKKOJAAecIBlhAbRADRwAEry6Bo+JPrRgN4EAJ3mll4cJlzw+qZwAOtm5spq+tn+gSEEDQDknRIQMnIKIgTiAFrl3AbVdRaNcADeBoYphg0AvkPipOOVU/WzZYBURHAjwu4gIEhk0pKOWAHoMmVKZHCAYURwtlQ3cpyzWLQ4OBwXwwKAQOCxVDRQEfJCSAiSOBQuAAd2AaBhAB9gURFFCyioAHxwAjRJC2eA3LDwmD3PpBDz4rE49hQsZEmFAoFQFFwAHabIAupyue54OJnoBUQjggCTCMy1eXhSIkEhEgz6EgAKjSGQ0YGA5kKy0sXK5qCQoSRErIxH4khg4PEKHIYAMgFYNwC7O4tjQ1wv8Raa5KJfB49WD+AgANbSIioO1wAD80jgrCIdodARs9pjcckAa5pPJ8FT6Yg0l8wWkzvw0gAXmwIBp7OB86b9EwWGwNKRbfbG9WwNIQI50JmiPaeukFHrgNIpCnexmiCgAMJw2kDqswF0m01Ii2M0Q9Z6AACJAAW4cBIvjgA5S4hIgAgiX5+gF7rnBztQaSOMAopACRE+QIUIiBgRwFHQWhpGoDdt3wOBDEFVF0VQS8kSIVsSTJCk4EAC/INTgbUp0yfVDQKJJMKBCs5z6DRwHYfRh1Hed9AAORQJt1xoyQ6LAVkiAaYiZ248wKLfU1/m8UBXT5UQUhaAIglCFJqIrFIhOyGMUCHEc52AQVAEvyfNCxwkACBAOlvn8CBRGkEs+ytK1OmENACGgcyrmAGR7CIEhni0SQVxsPwKQIfgXlUGEAqC3yCFC/gsluGlLIZAwvJikLWH4aQdGADRQCQagPAUtpLQAHkxEUfJgXxHAgRwZGABLqtq+rJB4IgwwgCNIzgCqRSElqiDqhrrl8JsHEHIaRvakh+tNGKauGtqOvyohgkAAIIPhsJaZvLZqdtahqOpFfRlTgZNkTVPxpBs6D4ChYoDEklxnDLGkGlIa5bnuR4gA}{link to Lean live}]

Using this bridge from global to local continuity, we can
understand the following definitions:
\begin{lstlisting}[language=lean]
def ContinuousAt (f : X → Y) (x : X) :=
  Tendsto f (𝓝 x) (𝓝 (f x))

theorem Continuous.continuousAt (h : Continuous f) :
    ContinuousAt f x :=
  h.tendsto x
\end{lstlisting}
The remaining local continuity concepts are defined similarly:
\begin{lstlisting}[language=lean]
def ContinuousWithinAt (f : X → Y) (s : Set X) (x : X) : Prop :=
  Tendsto f (𝓝[s] x) (𝓝 (f x))

def ContinuousOn (f : X → Y) (s : Set X) : Prop :=
  ∀ x ∈ s, ContinuousWithinAt f s x
\end{lstlisting}
Note that \lstinline[language=lean]|𝓝[s] x| denotes the
neighborhood filter of $x$ restricted to the set $s$, allowing us
to study the behavior of functions on arbitrary subsets.
We will use \lstinline[language=lean]|ContinuousOn| in the next
chapter.

% Section transition
Having introduced the topological machinery available in Mathlib, we
now turn to the specific topological concepts needed for our
formalization: connectedness and path-connectedness.

\section{Connectedness and
  Path-Connectedness}\label{sec:connectivity}

In this section, we recall the definitions of connectedness and
path-connectedness from standard topology, following
Munkres~\cite{munkres}, and present how they are formalized in
Mathlib.

\begin{definition}\label{def:connected}
  A topological space $X$ is \textbf{connected} if it cannot be
  written as the union of two disjoint nonempty open sets.
  Equivalently, the only subsets of $X$ that are both open and
  closed are $\varnothing$ and $X$ itself.
\end{definition}

\begin{definition}\label{def:path_connected}
  A topological space $X$ is \textbf{path-connected} if for every
  pair of points $a, b \in X$, there exists a continuous map
  $\gamma : [0,1] \to X$ (called a \textbf{path}) such that
  $\gamma(0) = a$ and $\gamma(1) = b$.
\end{definition}

\begin{figure}[H]
  \centering
  \begin{minipage}[t]{0.45\linewidth}
    \centering
    \textbf{\textcolor{UniBlue}{Connected}}\\[0.5em]
    \begin{tikzpicture}[scale=1.5]
      \draw[UniBlue, line width=1.5pt, fill=CurveBlue!20]
      plot[smooth cycle, tension=0.7]
      coordinates {(0.0, 0.5) (0.4, 1.0) (1.0, 0.75)
          (1.6, 1.0) (2.0, 0.5) (1.6, 0.0) (1.0, 0.25)
          (0.4, 0.0)};
      \node[font=\normalsize] at (1.0, 0.5) {$X$};
    \end{tikzpicture}
  \end{minipage}
  \hfill
  \begin{minipage}[t]{0.45\linewidth}
    \centering
    \textbf{\textcolor{UniBlue}{Path-Connected}}\\[0.5em]
    \begin{tikzpicture}[scale=1.5]
      \draw[UniBlue, line width=1.5pt, fill=CurveBlue!20]
      plot[smooth cycle, tension=0.7]
      coordinates {(0.0, 0.5) (0.4, 1.0) (1.0, 0.75)
          (1.6, 1.0) (2.0, 0.5) (1.6, 0.0) (1.0, 0.25)
          (0.4, 0.0)};
      \draw[dotred, line width=1.5pt]
      plot[smooth, tension=0.8]
      coordinates {(0.4, 0.5) (0.7, 0.6) (1.3, 0.4)
          (1.6, 0.5)};
      \fill[dotred] (0.4, 0.5) circle (1.5pt)
      node[left=1pt, font=\normalsize] {$a$};
      \fill[dotred] (1.6, 0.5) circle (1.5pt)
      node[right=1pt, font=\normalsize] {$b$};
      \node[dotred, font=\normalsize, above] at (0.5, 0.5)
      {$\gamma$};
    \end{tikzpicture}
  \end{minipage}
  \caption{Visual representation of connectedness vs.\
    path-connectedness.}
  \label{fig:connectivity-diagrams}
\end{figure}

It is a standard result that path-connectedness implies
connectedness~\cite[Theorem~25.5]{munkres}.
Intuitively, if a path-connected space could be split into two
disjoint open sets $U$ and $V$, then any path from a point in $U$
to a point in $V$ would yield a continuous map from the connected
space $[0,1]$ onto a disconnected space, which is impossible.
However, the converse does not hold: a space can be connected
without being path-connected.
The classical counterexample demonstrating this asymmetry is the
\textbf{topologist's sine curve}, which we formalize in the next
chapter.

\subsection{Mathlib's formalization}

In Mathlib, a subset $S$ of a topological space is connected if it
is nonempty and \textbf{preconnected}, where preconnected means it
cannot be covered by two disjoint open sets that each meet $S$:
\begin{lstlisting}[language=lean]
def IsPreconnected (s : Set α) : Prop :=
  ∀ (u v : Set α), IsOpen u → IsOpen v → s ⊆ u ∪ v →
    (s ∩ u).Nonempty → (s ∩ v).Nonempty → (s ∩ (u ∩ v)).Nonempty

structure IsConnected (s : Set α) : Prop where
  nonempty : s.Nonempty
  isPreconnected : IsPreconnected s
\end{lstlisting}

Path-connectedness is defined via the
\lstinline[language=lean]|IsPathConnected| predicate.
A subset $S$ is path-connected if there exists a point $x \in S$
such that every other point $y \in S$ can be joined to $x$ by a
path lying entirely in $S$:
\begin{lstlisting}[language=lean]
def IsPathConnected (F : Set X) : Prop :=
  ∃ x ∈ F, ∀ ⦃y⦄, y ∈ F → JoinedIn F x y
\end{lstlisting}
The auxiliary predicate \lstinline[language=lean]|JoinedIn| requires
the existence of a path between the two points that stays within
$S$:
\begin{lstlisting}[language=lean]
def JoinedIn (S : Set X) (x y : X) : Prop :=
  ∃ γ : Path x y, ∀ t, γ t ∈ S
\end{lstlisting}
Here, \lstinline[language=lean]|Path x y| denotes a continuous map
$\gamma : [0,1] \to X$ with $\gamma(0) = x$ and $\gamma(1) = y$.
The domain $[0,1]$ is represented in Lean by the
\lstinline[language=lean]|unitInterval| \textbf{subtype} of
\lstinline[language=lean]|ℝ|, defined as
$\{t : \mathbb{R} \mid 0 \leq t \leq 1\}$.
As we will see in the next chapter, working with this subtype
introduces some friction, since it does not automatically inherit
all arithmetic operations from $\mathbb{R}$.

\subsection{Useful theorems about connectedness}

Two results from Mathlib will be especially important for our
formalization.
The first concerns continuous images:

\begin{theorem}\label{thm:connected_image}
  The continuous image of a connected set is connected.
\end{theorem}

In Mathlib, this is available as
\lstinline[language=lean]|IsConnected.image|.
As a consequence, \lstinline[language=lean]|isConnected_Ioi.image|
states that the image of an interval $(a, \infty)$ under a
continuous map is connected.

The second result concerns closures:

\begin{theorem}\label{thm:connected_closure}
  Let $C$ be a connected subset of a topological space.
  If $C \subseteq S \subseteq \overline{C}$, then $S$ is connected.
\end{theorem}

In Mathlib, this is
\lstinline[language=lean]|IsConnected.subset_closure|.
Intuitively, adding limit points to a connected set cannot
disconnect it, because any open set that meets the added points
must also meet the original set.

% Chapter closing paragraph
In this chapter, we have surveyed how Lean and Mathlib formalize
topology: from the \lstinline[language=lean]|TopologicalSpace| type
class and the definition of continuity, through filters as a uniform
language for convergence, to the definitions of connectedness and
path-connectedness.
In the next chapter, we put these tools to work by formalizing the
topologist's sine curve---a space that is connected but not
path-connected.